\begin{mistake}{2026-04-19}{10}

\begin{problemcontent}
设 \(A\) 是 \(n\) 阶矩阵，满足 \(A^2 = A\)，求 \((A-2E)^3 - 3A\) 的值。
\end{problemcontent}

\begin{solutioncontent}
因为矩阵 \(A\) 与单位矩阵 \(E\) 满足交换律（即 \(AE = EA\)），可以直接利用二项式公式展开：
\[
\begin{aligned}
(A-2E)^3 &= A^3 - 3A^2(2E) + 3A(2E)^2 - (2E)^3 \\
&= A^3 - 6A^2 + 12A - 8E
\end{aligned}
\]

由已知条件 \(A^2 = A\)，可递推得到 \(A^3 = A^2 \cdot A = A \cdot A = A^2 = A\)。

将高次项降幂代入展开式中化简：
\[
\begin{aligned}
(A-2E)^3 &= A - 6A + 12A - 8E \\
&= 7A - 8E
\end{aligned}
\]

代入原题求值表达式：
\[
\begin{aligned}
(A-2E)^3 - 3A &= (7A - 8E) - 3A \\
&= 4A - 8E
\end{aligned}
\]
\end{solutioncontent}

\mistakeknowledge{
\begin{itemize}
 \item \textbf{核心公式/定理}：多项式展开定理（前提：两矩阵可交换 \(AB=BA\)）与幂等矩阵性质（\(A^k=A\)）。
 \item \textbf{易错点}：展开常数项时极易漏写单位矩阵 \(E\)，将 \((2E)^3\) 错写成纯常数 \(8\)，导致矩阵与标量非法运算。
 \item \textbf{关联知识}：矩阵乘法一般不满足交换律。常见的可交换特例包括：\(A\) 与 \(E\)、\(A\) 与 \(A^{-1}\)、\(A\) 与 \(A^*\) 以及 \(A\) 的两个多项式。
\end{itemize}}

\end{mistake}
